\documentclass{article}
\usepackage[utf8]{inputenc}
\usepackage{cite}
\usepackage{hyperref}
\usepackage{graphicx}

% --- Title ---
\title{The Role of Artificial Intelligence Chatbots in Enhancing Foreign Language Learning: Opportunities and Challenges}
\author{Alimzhan Seilkhan}
\date{April 2026}

\begin{document}

\maketitle

% --- 2. Abstract ---
\begin{abstract}
In the context of Industry 4.0, the application of advanced information and computer technology in education has become crucial[cite: 363, 371]. This paper examines the capabilities and various types of AI-based chatbots used as language learning media[cite: 366]. AI tools offer significant advantages by providing personalized learning journeys, dynamic content, and instantaneous feedback mechanisms that enhance language acquisition[cite: 554, 571]. Research indicates that chatbots increase learner confidence and provide 24/7 opportunities for language practice[cite: 368, 369]. Finally, this article discusses the ethical considerations and emerging future trends in the integration of AI within language instruction[cite: 559].
\end{abstract}

% --- 3. Introduction ---
\section{Introduction}

The rapid development of internet and computer technologies over the past three decades has fundamentally transformed the educational landscape[cite: 565]. Specifically, advancements in Artificial Intelligence (AI) and Natural Language Processing (NLP) have led to the emergence of Intelligent Computer Assisted Language Learning (ICALL) as a significant subfield. These innovations allow for more adaptive and individualized learning environments compared to traditional one-size-fits-all approaches[cite: 575].

A chatbot is an AI-based computer program designed to conduct conversations via text or audio[cite: 365, 377]. In language education, chatbots serve as virtual tutors available 24/7, assisting learners in identifying errors, providing feedback, and enhancing macro-skills such as speaking and writing[cite: 8, 142]. By simulating natural dialogue scenarios, these tools help learners improve their vocabulary, grammar, and overall communicative competence[cite: 180, 571].

Despite the immense potential of AI, its implementation in education faces several challenges. These include concerns regarding data privacy, the potential lack of cultural nuance in AI-generated texts, and the need for better teacher preparation to integrate these tools effectively[cite: 28, 555, 591]. The objective of this research is to analyze the effectiveness of AI-based language learning tools and to propose optimal methods for their pedagogical application.

\section{Related Works}

The classification of AI-based chatbots is essential to understanding their pedagogical roles. According to Haristiani (2019), chatbots can be categorized into three main types based on their structure and functionality: Flow chatbots, AI chatbots, and Hybrid chatbots[cite: 403, 406]. 

Flow chatbots follow a predefined tree-based structure where the developer sets fixed responses, often using buttons and keywords. In contrast, AI chatbots utilize machine learning to update their knowledge from previous interactions, allowing for more free-flowing engagement. Hybrid models combine these concepts, providing a balance between structured paths and intelligent communication.

Table 1 summarizes these categories as defined in the literature:

\begin{table}[h!]
\centering
\caption{Classification of Chatbots in Education}
\begin{tabular}{|l|l|p{6cm}|}
\hline
\textbf{Category} & \textbf{Sub-category} & \textbf{Primary Function} \\ \hline
Structure & Flow Chatbot & Tree-based with fixed responses; uses buttons/keywords. \\ \hline
Structure & AI Chatbot & Learns from experience; allows free engagement. \\ \hline
Purpose & Specialist & Focuses on a specific task (e.g., customer service or language learning). \\ \hline
Audience & Generalist & Provides general knowledge (e.g., Siri, Cortana). \\ \hline
\end{tabular}
\end{table}

Recent studies highlight that while most educational chatbots focus on English, there is an emerging trend in developing tools for other languages, such as "Gengobot" for Japanese grammar[cite: 427, 451, 455]. Furthermore, research emphasizes that chatbots are often perceived as more patient than human tutors, as they are willing to repeat materials endlessly without social pressure[cite: 485, 486].

% --- References ---
\begin{thebibliography}{99}
\bibitem{Haristiani2019} Haristiani, N. (2019). Artificial Intelligence (AI) Chatbot as Language Learning Medium: An inquiry. \textit{Journal of Physics: Conference Series}.
\bibitem{Creely2024} Creely, E. (2024). Exploring the Role of Generative AI in Enhancing Language Learning: Opportunities and Challenges. \textit{International Journal of Changes in Education}.
\bibitem{Woo2021} Woo, J. H., \& Choi, H. (2021). Systematic Review for AI-based Language Learning Tools. \textit{arXiv Preprint}.
\end{thebibliography}

\end{document}